\documentclass[10pt]{article}
\usepackage[margin=12mm]{geometry}
\usepackage{graphicx}
\usepackage{caption}
\pagestyle{empty}
\captionsetup{font=small,labelfont=bf}

\begin{document}
\begin{figure}[p]
\centering
\includegraphics[width=\textwidth]{fig8_ce_full_tree_collective_panel.pdf}
\caption{\textbf{Alternative full-tree reconstructions define a collective
sampled Pareto envelope.} Six reconstruction strategies were evaluated for the
same 1,000-edge \emph{C.\@ elegans} protein lineage: bottom-up optimization by
lineage layer, independent layerwise parent assignment, paired bottom-up
rebuilding, a degree-constrained spanning forest, top-down rebuilding, and a
terminal-only rebuild without reuse of measured internal protein states. Thin
colored curves retain the non-dominated solutions from each strategy's sampled
weight sweep. Colored curves retain the identity of each reconstruction
strategy. Costs are reported
in first-cousin-shuffle standard deviations and translated so the natural
lineage (black cross) is at the origin. The first-cousin null and fixed-topology
parametric Brownian reference are shown on the main axes; the more distant
internal-layer shuffle, full assignment shuffle, and random rebuild appear in
the inset. The Brownian reference fits a full 23-dimensional rate covariance
from all 1,000 measured branch-standardized increments, fixes four root states,
and simulates all scored descendants down the same topology. It is included
only as a geometric probe of proximity to the sampled Pareto envelope; its
zero-drift, common-rate diffusion assumptions are not intended as a realistic
model of embryogenesis. This replaces an incompatible historical 10-PC cache
and independent-node ancestral-state sampler. The
historical notebook's MST trace was also not comparable because it omitted
internal spanning-tree edges; its publication replacement enforces four roots,
at most two children per parent, and exactly 1,000 scored edges. The combined
non-dominated envelope of the colored strategies is dominated by the
degree-constrained spanning forest, top-down rebuild, and terminal-only
rebuild.}
\label{fig:ce_full_tree_collective}
\end{figure}
\end{document}
